\honest{Collapse Postulates}{Eliezer Yudkowsky}{2008}

I begin with a definition: ``Macroscopic decoherence'' is ``also known as'' many-worlds, the idea that the microscopic laws hold at every level. \nb{In the literature, decoherence is a process that every interpretation accepts, and many-worlds is one interpretation.} Before decoherence, I say, it had not occurred to anyone that those laws ``might apply universally.'' \nb{Von Neumann applied them to the measuring apparatus in 1932. The result, a superposition of macroscopically different states, is what the collapse postulate was added to avoid.}

The first idea of collapse, I suggest, went like this: the squared modulus of the amplitude is the probability, and once you know something did not happen, its probability goes to zero. ``Read literally,'' that makes knowledge, ``or even conscious awareness,'' the cause of collapse, and this ``was in fact the form of the theory put forth by Werner Heisenberg!'' \nb{Heisenberg linked the jump in the probability function to a change of knowledge, but wrote in 1958 that the transition from possible to actual ``is not connected with the act of registration of the result by the mind of the observer.''} Later, nervous about dualism, physicists made collapse an objective process.

Why only one surviving world? Because, I say, collapse theories ``were devised in a time when it simply didn't occur to any physicists that more than one world could exist!'' \nb{That fits the textbook postulate of 1932. The collapse theories studied today, Pearle's from 1976 and Ghirardi, Rimini and Weber's from 1986, came after Everett's paper of 1957.}

Experiments, I say, keep finding superposition in larger systems. So why does no one propose that collapse happens only at planetary scale, once a minute? \nb{Such a theory would not give single outcomes during the minute, and single outcomes are what collapse theories are for, as the post itself says a few paragraphs earlier.} A cynic might say collapse survives because people do not remember splitting. Then I call that ad hominem, and ``too early.''

Does any collapse theory have experimental support? ``No.'' \nb{Nor does any experiment favor many-worlds over standard quantum mechanics. Collapse theories, unlike many-worlds, predict effects that experiments are now looking for.}

Then my list: collapse would be the only non-linear, non-unitary, discontinuous, configuration-nonlocal, CPT-violating, Liouville-violating, random and faster-than-light phenomenon in physics. \nb{Most of this describes objective collapse accurately, and its proponents agree. But the continuous version of collapse is not discontinuous; a time-symmetric version has been proposed; and every theory in which measurements have single outcomes is nonlocal, so collapse is the only nonlocal phenomenon only on views, such as many-worlds, that escape Bell's theorem.} ``What does the god-damned collapse postulate have to do for physicists to reject it? Kill a god-damned puppy?'' \nb{Collapse theorists call the textbook postulate ``absolutely unacceptable'' too. The costs of many-worlds come two days later, in ``Many Worlds, One Best Guess.''}
